\documentclass[11pt,a4paper]{article}

\usepackage[french]{babel}
\usepackage[utf8]{inputenc}
\usepackage[T1]{fontenc}
\usepackage{lmodern}
\usepackage[margin=2.5cm]{geometry}
\usepackage{amsmath}
\usepackage{booktabs}
\usepackage{longtable}
\usepackage{array}
\usepackage{xcolor}
\usepackage{hyperref}
\usepackage{fancyhdr}
\usepackage{enumitem}

\definecolor{solidateal}{RGB}{15,76,79}
\definecolor{gris}{RGB}{90,90,90}
\definecolor{vertaccord}{RGB}{45,110,80}
\definecolor{orangevigilance}{RGB}{180,120,20}
\definecolor{rougerefus}{RGB}{160,50,45}

\hypersetup{
  colorlinks=true,
  linkcolor=solidateal,
  urlcolor=solidateal,
  pdftitle={SOLIDA -- Pourquoi ces seuils, et d'où viennent-ils},
  pdfauthor={SOLIDA}
}

\pagestyle{fancy}
\fancyhf{}
\fancyhead[L]{\small\textcolor{gris}{SOLIDA -- Pourquoi ces seuils, et d'où viennent-ils}}
\fancyhead[R]{\small\textcolor{gris}{\thepage}}
\renewcommand{\headrulewidth}{0.4pt}

\newcommand{\fichier}[1]{\noindent\textbf{Implémentation :} \\ \tiny\ttfamily#1\normalfont\normalsize\par}
\newcolumntype{L}[1]{>{\raggedright\arraybackslash}p{#1}}

\newcommand{\tranche}[2]{\textcolor{#1}{\textbf{#2}}}

\title{\textcolor{solidateal}{\Huge SOLIDA} \\[0.3em] \Large Pourquoi ces seuils de décision, et d'où viennent-ils}
\author{}
\date{Dernière mise à jour : \today}

\begin{document}

\maketitle

\section*{Objet de ce document}

Ce document répond à quatre questions sur le découpage en tranches de décision
(\tranche{vertaccord}{ACCORD}, \tranche{orangevigilance}{ACCORD SOUS CONDITION},
\tranche{rougerefus}{REFUS}) : pourquoi ce découpage existe sous cette forme, d'où vient chaque
nombre qui le compose, comment la formule est utilisée concrètement dans le système, et si elle
est adaptée à la réalité d'une coopérative togolaise.
La réponse à la dernière question n'est pas uniformément positive -- ce document distingue
explicitement ce qui est \emph{économiquement fondé et robuste} de ce qui reste, aujourd'hui,
\emph{un point de départ de démonstration à calibrer}. Aucune valeur n'est présentée comme sourcée
si elle ne l'est pas réellement dans le code ou la documentation de référence du projet.

\tableofcontents

\section{Le problème que la formule résout}

\subsection{Un seuil à 50\% n'a pas de sens ici}

Le modèle produit une probabilité de défaut $p$ pour chaque dossier. La tentation la plus simple
serait de décider : « accord si $p < 0{,}50$ ». C'est le comportement par défaut de n'importe quel
classifieur binaire, et c'est précisément le choix qu'il faut éviter dans un contexte de crédit.

Deux raisons s'y opposent :

\begin{enumerate}
\item \textbf{Le problème est déséquilibré.} Sur les données du générateur, le taux de défaut est
de l'ordre de 7\,\%. Un seuil à 0,50 suppose implicitement que la classe positive et la classe
négative sont également probables et également coûteuses -- aucune des deux hypothèses n'est
vraie ici.
\item \textbf{Les deux types d'erreur n'ont pas le même coût.} Accorder à un futur défaillant
(faux négatif) fait perdre du capital. Refuser un bon client (faux positif) fait perdre une marge
d'intérêt. Ce sont deux grandeurs économiques de nature différente, et rien ne dit qu'elles sont
égales.
\end{enumerate}

La bonne question n'est donc pas « le modèle pense-t-il que ce client va faire défaut ? » mais
« économiquement, à partir de quelle probabilité est-il plus coûteux d'accorder que de
refuser ? ». C'est une question de \textbf{théorie de la décision statistique appliquée au
risque de crédit}, pas une question de classification pure -- le cadre est connu sous le nom de
\emph{classification sensible au coût} (\emph{cost-sensitive classification}), formalisé
notamment par Elkan (2001).

\section{La matrice de coûts : d'où vient la forme du calcul}

\subsection{Les deux erreurs possibles}

\begin{center}
\begin{tabular}{@{}L{4.2cm}L{5cm}L{5cm}@{}}
\toprule
\textbf{Erreur} & \textbf{Ce qui se passe} & \textbf{Coût} \\
\midrule
Faux négatif ($C_{FN}$) & On accorde à un futur défaillant & Perte en capital = LGD $\times$ exposition \\
Faux positif ($C_{FP}$) & On refuse un bon client & Marge nette perdue sur la vie du prêt \\
\bottomrule
\end{tabular}
\end{center}

\subsection{Deux raffinements nécessaires}

\textbf{Raffinement 1 -- la perte en cas de défaut n'est pas 100\,\% du capital.} Quand un
sociétaire fait défaut, la coopérative ne perd pas la totalité du montant prêté : la caution
solidaire du groupe est appelée, une partie des échéances a déjà été remboursée, et il existe une
procédure de recouvrement. On modélise cette perte réelle par un taux \textbf{LGD}
(\emph{Loss Given Default}), réaliste entre 60 et 80\,\%, jamais 100\,\%. Symétriquement, le coût
d'un faux positif n'est pas l'intérêt brut du prêt mais la \textbf{marge nette} -- l'intérêt moins
le coût de refinancement de la coopérative elle-même.

\textbf{Raffinement 2 -- la calibration est un prérequis, pas un détail.} Un seuil sur $p$ n'a de
sens que si $p$ reflète un vrai risque. Un modèle qui annonce 0,17 alors que le risque réel est de
0,30 produit un seuil économiquement faux quel que soit le raffinement appliqué ensuite. C'est
pourquoi la calibration du modèle (courbe de calibration en déciles, correction isotonique si
nécessaire) est une étape bloquante avant toute mise en production, documentée séparément dans
\texttt{03-MODELE/03-scorecard-et-grille.md}.

\subsection{Dérivation du seuil optimal}

On refuse un dossier dès que le coût espéré du refus devient inférieur au coût espéré de l'accord :

\begin{equation}
(1 - p) \cdot C_{FP} < p \cdot C_{FN}
\end{equation}

En isolant $p$, on obtient le seuil de bascule :

\begin{equation}
\text{seuil}^{*} = \frac{C_{FP}}{C_{FP} + C_{FN}} = \frac{\text{marge}}{\text{marge} + \text{LGD}}
\end{equation}

Ce seuil est une \textbf{probabilité}, pas un score. Il se traduit ensuite en score par la même
transformation PDO (\emph{Points to Double the Odds}) utilisée pour l'ensemble de la scorecard, ce
qui permet de l'afficher et de le comparer directement aux autres seuils du système.

\fichier{backend/solida/domain/rules/grille.py -- \texttt{ParametresGrille.seuil\_economique()}}

\subsection{Des trois tranches, pas d'une seule bascule}

Une seule bascule accord/refus serait trop rigide pour un métier où l'agent doit pouvoir
distinguer un dossier clairement favorable d'un dossier qui l'est mais mérite une vigilance
particulière. Le système multiplie donc le seuil économique par un coefficient qui délimite
trois zones :

\begin{equation}
p < \text{seuil}^{*} \times m
\quad\text{avec}\quad
m \in \{\,m_{\text{accord}},\ 1\,\}
\end{equation}

\begin{center}
\begin{tabular}{@{}lll@{}}
\toprule
\textbf{Tranche} & \textbf{Condition} & \textbf{Multiplicateur} \\
\midrule
\tranche{vertaccord}{ACCORD} & $p < \text{seuil}^{*} \times 0{,}6$ & $m_{\text{accord}} = 0{,}6$ \\
\tranche{orangevigilance}{ACCORD SOUS CONDITION} & $p < \text{seuil}^{*}$ & implicite (= 1) \\
\tranche{rougerefus}{REFUS} & sinon & -- \\
\bottomrule
\end{tabular}
\end{center}

La lecture économique est directe : en dessous de 60\,\% du seuil de bascule, le dossier est
suffisamment loin du point où l'accord devient risqué pour être recommandé sans réserve. Entre
$0{,}6\times$ le seuil et le seuil lui-même, le dossier est encore favorable mais assez proche du
point de bascule pour justifier une condition (l'agent tranche avec un levier identifié). Au-delà
du seuil économique pur -- le point exact où, selon la matrice de coûts, refuser coûte moins cher
qu'accorder --, le signal de risque est jugé trop fort pour une recommandation favorable, avec ou
sans condition.

\textbf{Une quatrième zone existait, elle a été supprimée.} Une version antérieure de la grille
ajoutait une zone \tranche{orangevigilance}{COMITÉ DE CRÉDIT} entre $1{,}0\times$ et $1{,}6\times$
le seuil, présentée comme la seule renvoyant à un examen approfondi. Elle a été retirée : le
comité de crédit valide en réalité \textbf{chaque} tranche (section~\ref{sec:humaine}), ce n'était
donc pas une zone de décision à part mais une confusion de vocabulaire. L'ancienne zone d'examen
est désormais absorbée par \tranche{rougerefus}{REFUS} -- un dossier qui s'y trouve garde
exactement le même parcours qu'avant (réexamen possible, voir plus loin), seul le nom de l'étape
change.

\fichier{backend/solida/domain/rules/grille.py -- \texttt{decider()}}

\subsection{La décision reste humaine à chaque tranche}
\label{sec:humaine}

Aucune tranche, y compris \tranche{vertaccord}{ACCORD}, ne débloque un crédit automatiquement.
Le moteur produit une \textbf{recommandation}, jamais une décision exécutoire : l'agent, puis le
comité de crédit, reste seul décideur final pour les trois tranches indifféremment, avec la fiche
de justification comme support. C'est précisément pourquoi une tranche nommée d'après le comité
de crédit n'avait pas lieu d'être : elle laissait croire que le comité n'intervenait que là, alors
qu'il valide chaque dossier.

\section{D'où viennent les nombres : ce qui est sourcé, ce qui ne l'est pas}

C'est la question la plus importante de ce document, et celle où l'honnêteté prime sur la
présentation. Le code lui-même le dit noir sur blanc :

\begin{quote}\itshape
« Mais les seuils ne sont pas un choix technique. \texttt{marge}, \texttt{LGD} et le
multiplicateur de zone d'accord (0,6) traduisent l'arbitrage entre taux d'approbation et risque
accepté, qui appartient à la coopérative, pas au code. Les valeurs de \texttt{decision.py}
(\texttt{marge = 0,15}, \texttt{LGD = 0,75}) sont un point de départ pour construire et tester la
mécanique -- pas la vérité finale. »
\end{quote}
\hfill --- \texttt{03-MODELE/03-scorecard-et-grille.md}

\subsection{Ce qui est réellement démontré}

\begin{itemize}
\item La \textbf{forme} du calcul -- probabilité $\to$ coût espéré $\to$ seuil optimal $\to$
traduction en score $\to$ trois tranches par multiplicateur croissant -- est un résultat de
théorie de la décision statistique, indépendant de tout chiffre togolais ou uémoaïen. Sa validité
ne dépend pas des valeurs numériques choisies.
\item L'\textbf{ordre de grandeur} de LGD (60 à 80\,\%, contre 100\,\% naïvement supposé) est
défendable économiquement dès lors qu'une caution solidaire et un recouvrement partiel existent
réellement dans le produit -- ce qui est le cas de SOLIDA (couche solidaire, groupe de caution).
\item Sur le portefeuille de test du générateur (1\,885 crédits synthétiques, 173\,M FCFA), passer
du seuil naïf 0,50 au seuil optimal $0{,}15/(0{,}15+0{,}75) = 0{,}167$ ferait économiser environ
2,4\,M FCFA, soit 17\,\% du coût total des erreurs sur cet échantillon. C'est un résultat de
simulation sur données synthétiques, utile pour démontrer le \emph{mécanisme}, pas une mesure
d'impact sur le portefeuille réel de la coopérative.
\item Le seuil théorique (0,167) et le seuil optimal recalculé empiriquement sur le jeu de
calibration (0,142) sont proches, ce qui suggère que la mécanique se comporte de façon cohérente
une fois appliquée à un modèle réellement entraîné -- mais ce rapprochement reste mesuré sur des
données synthétiques, pas sur l'historique réel d'une coopérative togolaise.
\end{itemize}

\subsection{Ce qui n'est pas sourcé, et doit être dit comme tel}

\texttt{marge = 15\,\%} et \texttt{LGD = 75\,\%} ne proviennent \textbf{d'aucune étude, d'aucun
historique de recouvrement, d'aucune donnée BCEAO ou togolaise}. Ce sont des valeurs de
démonstration choisies pour que la mécanique soit testable avant que la coopérative ne dispose
d'un chiffrage réel. Il en va de même pour le multiplicateur de zone d'accord (0,6) : il fixe la
largeur relative de la zone favorable sans réserve, mais rien dans \texttt{03-MODELE/} ne documente
pourquoi 0,6 plutôt que 0,5 ou 0,7 -- c'est un choix d'espacement raisonnable, non calibré.

\textbf{Ce que cela signifie concrètement} : le mécanisme (quelle formule, quelle traduction en
score, quel découpage en trois zones) est solide et transposable tel quel. Les \emph{curseurs}
(quelle marge exacte, quel LGD exact, quelle largeur de zone d'accord) ne le sont pas -- ils
demandent un arbitrage réel de la coopérative avant d'être présentés comme définitifs.

\subsection{Un contre-exemple instructif : le ratio d'endettement, lui, est sourcé}

Le système contient un autre seuil, structurellement proche (une fraction qui déclenche une
tranche de risque), mais dont le statut de sourcing est radicalement différent : le ratio maximal
mensualité / revenu, fixé à 33\,\%. Contrairement à \texttt{marge}/\texttt{LGD}, ce chiffre a été
retrouvé et documenté après coup dans le code :

\begin{quote}\itshape
« Sourcé (2026-09-14, trouvé après coup -- la valeur avait d'abord été posée en pure
démonstration) : étude de terrain CEF-MF Lomé 2025 [...] -- ratio mensualité/revenu moyen de 23\,\%
sur l'échantillon, rupture identifiée à 33\,\% au-delà de laquelle le risque d'impayé augmente. »
\end{quote}
\hfill --- \texttt{backend/solida/domain/rules/grille.py}, champ \texttt{ratio\_endettement\_maximal}

Ce contraste illustre exactement la différence à retenir : un seuil du même système, dans le même
fichier, peut être solidement ancré dans une réalité togolaise (33\,\%, étude CEF-MF) ou rester un
point de départ non calibré (marge, LGD, multiplicateur d'accord). Le premier peut être présenté au jury et
à la coopérative comme une donnée de terrain. Les seconds doivent être présentés comme un mécanisme
prêt à recevoir les vrais chiffres de la coopérative, pas comme des chiffres déjà validés.

\section{Comment la formule est utilisée concrètement}

\subsection{À l'exécution, pour chaque dossier}

\begin{enumerate}
\item Le modèle (EBM, SOCLE ou enrichi selon le profil) produit une probabilité de défaut
calibrée $p$.
\item \texttt{ParametresGrille.seuil\_economique()} calcule
$\text{seuil}^{*} = \text{marge}/(\text{marge}+\text{LGD})$ à partir de la configuration
\emph{active} en base (table \texttt{grille\_decision}), jamais d'une constante figée dans le code.
\item \texttt{decider()} compare $p$ au seuil économique et à sa version réduite par le
multiplicateur d'accord, et renvoie l'une des trois tranches.
\item La tranche, avec le score et la décomposition en points, alimente la fiche de justification
présentée à l'agent -- jamais une décision exécutoire directe.
\end{enumerate}

\subsection{Versionnage : rien n'écrase silencieusement une décision passée}

Toute modification de ces paramètres (\texttt{marge}, \texttt{LGD}, multiplicateur d'accord) crée
une nouvelle \texttt{version\_grille} en base, via l'endpoint d'administration
\texttt{/parametrage/grille}. L'ancienne version est désactivée, jamais supprimée, et chaque décision passée
conserve la référence de la version qui l'a produite. Sans cela, aucune décision historique ne
serait rejouable ni auditable -- un point non négociable pour un produit financier soumis à
justification devant un sociétaire ou un régulateur.

\fichier{backend/solida/adapters/persistence/grille\_repository\_sql.py}
\fichier{backend/solida/infrastructure/routers/parametrage.py}

\subsection{Qui peut ajuster ces chiffres, et comment}

Ces paramètres ne sont pas codés en dur : ils sont exposés par l'endpoint d'administration
\texttt{/parametrage/grille}, pensé pour un écran dédié où la coopérative visualise l'effet d'un
changement de \texttt{marge}/\texttt{LGD}/multiplicateur sur l'historique avant de l'activer. C'est
la réponse opérationnelle au constat de la section précédente : le mécanisme est prêt, l'écran
pour que la coopérative y injecte ses propres chiffres existe déjà, il reste à les déterminer.

\section{Pourquoi cette approche est pertinente pour SOLIDA}

\begin{itemize}
\item \textbf{Elle remplace un seuil arbitraire par un seuil justifiable.} Face à un comité de
crédit ou un auditeur, « le seuil est 0,167 parce que $C_{FP}/(C_{FP}+C_{FN})$ avec nos chiffres de
marge et de perte » est défendable ; « le seuil est 0,50 parce que c'est la valeur par défaut d'un
classifieur » ne l'est pas.
\item \textbf{Elle est cohérente avec un problème déséquilibré (7\,\% de défaut).} Un seuil fixe à
0,50 sur une population à 7\,\% de défaut classerait presque tout le monde en accord et n'apporterait
aucune valeur de tri -- le seuil économique, nettement plus bas (0,167), est ce qui rend le score
réellement discriminant sur cette population.
\item \textbf{Elle respecte le rôle du modèle dans SOLIDA : recommander, jamais décider.} Les
trois tranches, et en particulier la zone \tranche{orangevigilance}{ACCORD SOUS CONDITION},
formalisent explicitement un palier où le modèle signale une proximité du seuil de bascule et
renvoie un levier concret à l'agent plutôt que de forcer une décision binaire.
\item \textbf{Elle est réglable sans nouveau déploiement.} Le stockage en base versionné
(section~4.2) permet à la coopérative de faire évoluer son arbitrage risque/approbation au fil du
temps -- notamment quand les coûts réels changent (taux d'usure BCEAO, coût de refinancement) --
sans dépendre d'un cycle de développement.
\item \textbf{Elle prolonge le refus en parcours, pas en porte fermée.} Toute décision autre
qu'un \tranche{vertaccord}{ACCORD} sans réserve s'accompagne de conditions de réexamen concrètes
et vérifiables (régularité d'épargne, ratio d'épargne nantie, endettement à réduire) et d'une
trajectoire de plafond accessible en cas de remboursement sans incident.

\fichier{backend/solida/domain/rules/progressif.py -- \texttt{lister\_conditions\_reexamen}}
\end{itemize}

\section{Est-ce adapté à nos réalités ? Une réponse honnête}

\subsection{Ce qui est adapté}

Le \textbf{cadre} -- coût espéré, calibration préalable obligatoire, décision humaine finale,
traçabilité par version, ajustabilité sans redéploiement -- est directement adapté à une
coopérative financière ouest-africaine aux ressources informatiques limitées : aucune dépendance
lourde, aucune boîte noire, un mécanisme audité et rejouable. Le raffinement LGD (perte réelle
60-80\,\%, pas 100\,\%) est également adapté \emph{parce que} le produit dispose réellement d'une
caution solidaire et d'un recouvrement -- ce ne serait pas vrai pour un crédit sans garantie.

\subsection{Ce qui ne l'est pas encore}

\begin{enumerate}
\item \textbf{\texttt{marge = 15\,\%} et \texttt{LGD = 75\,\%} n'ont, à ce jour, aucune source
togolaise ou UEMOA.} Contrairement au ratio d'endettement (33\,\%, étude CEF-MF Lomé 2025), rien
dans \texttt{03-MODELE/} ne documente le coût de refinancement réel de la coopérative, ni son taux
de recouvrement réel après appel de la caution solidaire. Tant que ces deux chiffres restent des
valeurs de démonstration, le seuil de 0,167 est \emph{mécaniquement correct} mais
\emph{économiquement non calibré} sur la réalité de la coopérative.
\item \textbf{Le multiplicateur de zone d'accord n'est pas calibré sur un historique réel.} Il
fixe la largeur de la zone \tranche{orangevigilance}{ACCORD SOUS CONDITION} sans mesure de la
charge de travail que cela impose à l'agent -- un multiplicateur trop bas enverrait trop de
dossiers en condition pour la capacité réelle de traitement de la coopérative.
\item \textbf{Le gain chiffré (2,4\,M FCFA, 17\,\%) est un résultat de simulation}, pas une mesure
sur le portefeuille réel. Il démontre que le mécanisme fonctionne et vaut la peine d'être calibré ;
il ne doit pas être présenté comme un gain déjà obtenu par la coopérative.
\item \textbf{Aucun test n'a encore eu lieu à l'échelle réseau UEMOA} (5,7 millions de sociétaires
visés) : le mécanisme est validé sur un jeu de test synthétique de 1\,885 crédits, plusieurs ordres
de grandeur en dessous du volume réel visé.
\end{enumerate}

\subsection{Ce qu'il reste à faire pour que ce soit pleinement adapté}

\begin{itemize}
\item Faire chiffrer \texttt{marge} et \texttt{LGD} par la coopérative elle-même : coût de
refinancement réel, taux de recouvrement réel observé après appel de caution solidaire (pas une
hypothèse de 60-80\,\%).
\item Faire valider le multiplicateur de zone d'accord par le comité de crédit, en fonction de sa
capacité réelle à traiter un volume donné de dossiers en zone d'accord sous condition.
\item Recalculer le seuil optimal empirique (comme le 0,142 obtenu sur le jeu de calibration
synthétique) une fois le modèle entraîné sur des données réelles de la coopérative, et comparer à
la valeur théorique.
\item Documenter la source dans \texttt{03-MODELE/03-scorecard-et-grille.md} dès que ces chiffres
sont fixés -- exactement comme cela a été fait après coup pour le ratio d'endettement -- pour que
la formule cesse d'être un point de départ et devienne une décision tracée.
\end{itemize}

\section{Garde-fous}

Le seuil optimal reste un \textbf{point de départ recommandé}, jamais une porte automatique. Le
code et la documentation de référence posent trois garde-fous explicites, indépendants de la
question du calibrage des valeurs :

\begin{enumerate}
\item La décision reste humaine à chaque tranche (agent, puis comité selon la tranche).
\item Le seuil se surveille par segment pour éviter d'exclure systématiquement une zone
géographique ou un secteur d'activité -- cohérence avec l'engagement de non-discrimination du
projet (\texttt{03-MODELE/06-ethique-et-non-discrimination.md}).
\item Le seuil se recalcule quand les coûts réels de l'institution changent (taux d'usure, coût de
refinancement) -- ce n'est pas une valeur figée une fois pour toutes.
\end{enumerate}

\section*{Note de compatibilité historique}

Une version antérieure de ce document décrivait quatre tranches, avec une zone
\tranche{orangevigilance}{COMITÉ DE CRÉDIT} entre l'accord sous condition et le refus. Cette
zone a été retirée du mécanisme de décision (section~2.4) : le comité de crédit valide chaque
tranche, ce n'était pas une zone à part. La valeur \texttt{comite\_de\_credit} reste néanmoins
définie dans le code (\texttt{TrancheDecision}, backend) et dans les contrats d'API, pour une
seule raison : la table qui persiste les décisions (\texttt{decision\_scoring}) est en
\textbf{insertion seule} -- un trigger interdit toute modification ou suppression après coup, y
compris pour corriger une tranche devenue obsolète. Les décisions prises avant ce changement et
portant cette tranche restent donc lisibles telles quelles dans le registre et sur une fiche déjà
archivée ; aucune nouvelle décision ne peut plus la produire.

\section*{En résumé}

La \textbf{forme} du découpage en trois tranches -- une matrice de coûts, un seuil économique, sa
traduction en score, un multiplicateur qui ouvre un palier de vigilance avant le refus -- est un
mécanisme solide, explicable et directement pertinent pour un produit de microcrédit coopératif.
Les \textbf{valeurs} actuellement injectées dans ce mécanisme (\texttt{marge = 15\,\%},
\texttt{LGD = 75\,\%}, multiplicateur d'accord 0,6) sont, à la date de ce document, des valeurs de
démonstration non sourcées sur la réalité togolaise -- à la différence du ratio d'endettement de
33\,\%, qui l'est. Le mécanisme n'a donc pas besoin d'être reconstruit ; il a besoin d'être nourri
des vrais chiffres de la coopérative, par l'écran d'administration qui existe déjà à cet effet.

\end{document}
